\documentclass[a4paper, 12pt]{article}
\usepackage{a4wide}
\usepackage {amsmath}
\usepackage {graphicx}
\usepackage[utf8]{inputenc} 
\usepackage[francais]{babel}
\usepackage{fancyhdr}
\usepackage{setspace}
\usepackage{fancyhdr}
\usepackage{lastpage}
\usepackage{extramarks}
\usepackage{chngpage}
\usepackage{soul}
\usepackage[usenames,dvipsnames]{color}
\usepackage{graphicx,float,wrapfig}
\usepackage{ifthen}
\usepackage{listings}
\usepackage{courier}
\usepackage{esint}
\usepackage{bbm}
\usepackage{graphics}
\usepackage{graphicx}
\usepackage{subfig}
\usepackage{epsfig}
\usepackage{pgf,tikz}
\usetikzlibrary{arrows}
\usepackage{braket}
\usepackage{MnSymbol,wasysym}
\usepackage{marvosym}
\usepackage{dsfont}

\def\ChoixDeVersion{corrections} % Choix de la version a compiler (corrections ou pas)
\newcommand{\Corrections}[1]{%
        \color{red}
	\ifthenelse{
		\equal{\ChoixDeVersion}{nocorrections}
	}{CORRECTION\\ #1}{ }
        \color{black} 
}

\lhead{} 
\chead{} 
\rhead{\bfseries  $DM - \theta$-schéma \Corrections{}} 

\lfoot{J. C. Toussaint}
%\cfoot{} 
%\rfoot{\thepage}

\let\vaccent=\v % rename builtin command \v{} to \vaccent{}
\renewcommand{\v}[1]{\ensuremath{\mathbf{#1}}} % for vectors
\newcommand{\gv}[1]{\ensuremath{\mbox{\boldmath$ #1 $}}} 

\newcommand{\grad}[1]{\gv{\nabla} #1} % for curl
\let\divsymb=\div % rename builtin command \div to \divsymb
\renewcommand{\div}[1]{\gv{\nabla} \cdot #1} % for divergence
\newcommand{\curl}[1]{\gv{\nabla} \times #1} % for curl

% This is the color used for MATLAB comments below
\definecolor{MyDarkGreen}{rgb}{0.0,0.4,0.0}

% For faster processing, load Matlab syntax for listings
\lstloadlanguages{Matlab}%
\lstset{language=Matlab,                        % Use MATLAB
        frame=single,                           % Single frame around code
        basicstyle=\small\ttfamily,             % Use small true type font
        keywordstyle=[1]\color{Blue}\bf,        % MATLAB functions bold and blue
        keywordstyle=[2]\color{Purple},         % MATLAB function arguments purple
        keywordstyle=[3]\color{Blue}\underbar,  % User functions underlined and blue
        identifierstyle=,                       % Nothing special about identifiers
                                                % Comments small dark green courier
        commentstyle=\usefont{T1}{pcr}{m}{sl}\color{MyDarkGreen}\small,
        stringstyle=\color{Purple},             % Strings are purple
        showstringspaces=false,                 % Don't put marks in string spaces
        tabsize=5,                              % 5 spaces per tab
        %
        %%% Put standard MATLAB functions not included in the default
        %%% language here
        morekeywords={xlim,ylim,var,alpha,factorial,poissrnd,normpdf,normcdf},
        %
        %%% Put MATLAB function parameters here
        morekeywords=[2]{on, off, interp},
        %
        %%% Put user defined functions here
        morekeywords=[3]{FindESS, homework_example},
        %
        morecomment=[l][\color{Blue}]{...},     % Line continuation (...) like blue comment
        numbers=left,                           % Line numbers on left
        firstnumber=1,                          % Line numbers start with line 1
        numberstyle=\tiny\color{Blue},          % Line numbers are blue
        stepnumber=5                            % Line numbers go in steps of 5
        }

% Includes a MATLAB script.
% The first parameter is the label, which also is the name of the script
%   without the .m.
% The second parameter is the optional caption.
\newcommand{\matlabscript}[2]
  {\begin{itemize}\item[]\lstinputlisting[caption=#2,label=#1]{#1.m}\end{itemize}}

\pagestyle{fancy}

\begin{document}

\bibliographystyle{alpha}

\title{Intégrateur de type $\theta$-schéma pour les équations aux dérivées partielles}

\author{Jean-Christophe Toussaint\\
%  Phelma\\
%\texttt{jean-christophe.toussaint@phelma.grenoble-inp.fr}
}
\date{\today}
 

\maketitle

Ce document donne les éléments nécessaires pour construire un intégrateur de type $\theta$-schéma pour résoudre
l'équation de la chaleur dépendante du temps et pour étudier son stabilité dans le cadre d'une approche énergétique.

\section{Equation de la chaleur}

Le  $\theta$-schéma est très largement utilisé pour résoudre  des équations aux dérivées partielles
dépendantes du temps, comme l'équation de la chaleur.
Le second membre \eqref{edo1} est alors remplacé par un terme laplacien
qui se comporte comme un
terme de raideur.

\noindent On considère l'équateur de la chaleur dépendante du temps sans terme source:
\begin{equation}
\rho\ C_p \frac{\partial T}{\partial t}+ \div \gv{j}=0 \ \text{avec} \ \gv{j}=-k_{th} \grad{T}
\label{th1}
\end{equation}
où $\gv{j}$ est le vecteur flux de chaleur par unité de surface et $k_{th} $ dénote la conductibilité thermique, $\rho$ la masse volumique, $C_p$ la capacité calorifique.

En notant dans la suite $v=\frac{\partial T}{\partial t}$, la variation locale de température par unité de temps, l'équation \eqref{th1} devient:

\begin{equation}
\frac{\partial T}{\partial t}  = \frac{k_{th}}{\rho\; C_p } \; \Delta T 
\label{th2orig}
\end{equation}

\noindent ou encore en introduisant un paramètre $\lambda=\frac{k_{th}}{\rho\; C_p } $

\begin{equation}
v  =  \lambda \; \Delta T
\label{th2}
\end{equation}

\subsection{Le $\theta$-schéma:}

\noindent On propose par conséquent le $\theta-$schéma suivant :
 \begin{eqnarray}
 T^{n+1} = T^n+ k   v^{n+\theta} +
 \left\{ 
\begin{array}{l l l}
  \vartheta(k^3) \text{\, si \,} \theta=\frac{1}{2} \\
   \\
\vartheta(k^2) \text{\, sinon} \end{array} \right. 
  \label{schema}
  \end{eqnarray}
  
\noindent D'autre part, 
\begin{equation}
T \equiv T^{n + \theta}= T^n+ \theta k v^n + \vartheta(k^2)
 \label{DLord2orig}
\end{equation}

donne une estimation de $T$ à l'instant $n+\theta$. Sans modifier la précision
du développement,
on remplace $v^n$ dans \eqref{DLord2orig} par $v^{n+\theta}$. 

A l'instant $n+\theta$, l'équation \eqref{th2} devient:
\begin{equation}
v^{n + \theta}= \lambda\; \Delta (T^n+ \theta k v^{n + \theta}) + \vartheta(k^2)
 \label{DLord2inter}
\end{equation}

Ce qui revient à résoudre une équation en  $v^{n + \theta}$ en négligeant localement
les termes d'ordre 2 en temps. On applique la relation \eqref{schema} pour passer à 
l'instant suivant.

En résumé, le $\theta-$schéma comporte donc une étape de résolution :
\begin{equation}
 ( \mathds{1} -\lambda\; \theta\; k \Delta)\; v^{n + \theta} = \lambda\; \Delta T^n 
 \label{ThschemaStep1}
\end{equation}
puis une étape d'évolution :
\begin{equation}
T^{n+1} = T^n+ k   v^{n+\theta}
\label{ThschemaStep2}
\end{equation}

Comme le laplacien est un opérateur à valeurs propres négatives, on retrouve dans le terme
de gauche de l'équation \eqref{ThschemaStep1}, un préfacteur jouant le rôle de filtre passe-bas. Il coupe donc les hautes
fréquences spatiales à l'origine des instabilités des schémas numériques.

 \section{Application aux différences finies}
 
 Son application aux différences finies est immédiate. Après discrétisation de l'espace
 avec une grille régulière pour simplifier, on applique à chaque noeud $p \in [1, N]$ de la grille
 l'équation \eqref{ThschemaStep1} pour construire un système d'équations linéaires
 et le résoudre. On applique ensuite en chaque noeud l'équation d'évolution \eqref{ThschemaStep2}. 
 En résumé, on a:
 \begin{eqnarray}\left\{ 
\begin{array}{l l l}
   ( \mathds{1} -\lambda\; \theta\; k \Delta)\; v_p^{n + \theta} = \lambda\; (\Delta T)_p^n \\
   \\
 T_p^{n+1} = T_p^n+ k   v_p^{n+\theta} \end{array} \right. 
  \label{FDschema}
  \end{eqnarray}

\begin{enumerate} 
 \item On intéresse à l'évolution de la température au cours du temps dans une plaque.
 On suppose que l'approximation 1D est suffisante. Rappeler sous quelles conditions.
 \item On suppose que le système est périodique : on a  $T_{0}^n=T_N^n$ et $v_{0}^n=v_N^n$ à tout instant.
 Montrer comment l'expression discrète de l'opérateurs laplacien doit être adaptée
  pour tenir compte de la condition de périodicité.
 
 \item Développer un code simpliste mettant en oeuvre le $\theta$-schéma. Les paramètres de simulation
 seront regroupés dans une structure {\tt fd}.   Préciser la distribution initiale de température
 que vous avez choisie.
 On définira les fonctions :
 \begin{enumerate} 
 \item  {\tt lap1(N)} retournant la matrice associée à l'opérateur laplacien pour une grille de pas unitaire, tenant
 compte de la périodicité,
 \item {\tt solver(fd)} gérant la boucle sur le temps et l'affichage 
 de la distribution de température à chaque itération.
 \item {\tt fd\_main} programme principal appelant les fonctions décrites ci-dessus.
  \end{enumerate}

% L'utilisateur pourra appliquer des conditions de dirichlet ou de Neumann à gauche et à droite.
% {\it Noter que l'archive compressée au format ZIP du répertoire contenant le code source
% vous sera demandée lors du DS de méthodes numériques.}
 \item Résumer les résultats obtenus avec une de vos simulations que vous pensez pertinente 
 en précisant bien les paramètres de simulation et les conditions aux limites.
\end{enumerate} 

 \subsection{Stabilité énergétique du $\theta$-schéma et précision}

Dans [Euvrard], l'étude de stabilité énergétique en différences finies fait
appel à des astuces de majoration sur des expressions discrètes et est par conséquent
plus difficile à appréhender que celle au sens de Von Neumann.

On propose de développer une approche plus progressive en difficultés.
On se place d'abord dans la limité continue, on définit une co-énergie proportionnelle
à  $E \propto \int_0^L T(x, t)^2 \; dx$ pour un plaque homogène d'épaisseur $L$, infiniment haute.
{\it Pour simplifier, on suppose que des conditions périodiques sont appliquées
en $x=0$ et $x=L$.}

\begin{enumerate} 

\item Dans la limite continie, calculer la variation de la co-énergie avec le temps, $\frac{dE}{dt}$ et montrer qu'elle
est monotone décroissante.

\item Après discrétisation en différences finies, la température est évaluée sur une grille de pas $h$.
Les dérivées spatiales sont estimées à partir des développement de Taylor de la température.
Donner l'expression discrète de la variation $\Delta E = E_{n+1}-E_{n}$ avec une précision
en $\vartheta{(h^3)}$.

\item Pour éviter de traiter les noeuds extrémités comme des noeuds particuliers, on applique des conditions périodiques telles que $T_0 \equiv T_N$.

On montrera que la somme avec le laplacien discret s'écrit en fait, sous la forme 
$\sum_{j=1}^{N} (T_{j+1}^{n+\theta}-T_j^{n+\theta})^2$ où $L= N\, h$.

\item Pour $\theta \in ]\frac{1}{2}, \; 1]$, montrer que le schéma \eqref{FDschema} est dissipatif
 et donc inconditionnellement stable. Sa précision est d'ordre 1 en temps.
  
 \item Pour $\theta =\frac{1}{2}$, montrer qu'il est d'ordre 2 en temps. 
 Le schéma est dit semi-implicite.
 
\item Montrer que, pour $\theta \in [0,\;  \frac{1}{2}[$, le schéma est dissipatif à condition
de choisir le pas de temps $k<k_c$ où $k_c$ est une valeur critique dont on donnera
l'expression. Sa précision est d'ordre 1 en temps.
On utilisera l'inégalité suivante $(a-b)^2 \leq 2\; (a^2 + b^2)$ pour définir une expression de $k_c$.

Commenter et comparer la valeur de $k_c$ pour $\theta=0$ avec celle obtenue par une
analyse de stabilité de Von Neumann.

 \item  Implémenter le calcul de $\Delta E$ dans votre code de calcul.
 \item On veut vérifier numériquement la condition à imposer sur $\theta$
 pour observer la stabilité. Proposer un ensemble de simulations en reportant
 bien la liste des paramètres pertinents sous la forme d'un tableau.
 
 \item Le cas intéressant est $\theta = \frac{1}{2}$, 
seule cette valeur particulière de $\theta$ permet d'atteindre une précision est d'ordre $2$ en temps.
Proposez un protocole permettant de vérifier que le schéma est d'ordre 2 en temps.
 Appliquez-le en précisant les paramètres de simulation.

\end{enumerate}


% Une référence \cite{Press:2007}
 %\cite{Dhatt:2005}

%\bibliography{biblio}

\section{Conclusion}
           
Le  $\theta$-schéma est très largement utilisé pour résoudre  des équations différentielles
et les équations aux dérivées partielles
dépendantes du temps, comme les équations de diffusion. Il permet de cumuler précision
d'ordre 2 et stabilité inconditionnelle en imposant $\theta = \frac{1}{2}$. \\
 
\vspace{-0.5cm}
\section{Références}

\noindent Euvrard D. : {\it Résolution numérique des équations aux dérivées partielles de la physique, de la mécanique et des sciences de l'ingénieur : Différences finies, éléments finis, problèmes en domaines non bornés }, ed. Masson, 1994.

\end{document}

